\documentclass[11pt]{article}
\usepackage{amsmath,amssymb,mathtools,mathrsfs,geometry,hyperref}
\hypersetup{hypertexnames=false}
\geometry{margin=24mm}
\newcommand{\C}{\mathbb C}\newcommand{\diag}{\operatorname{diag}}
\newcommand{\Tr}{\operatorname{tr}}\newcommand{\Res}{\operatorname{Res}}
\title{The complete signed ES collision in both original conductor metrics}
\author{Independent programme derivation}\date{20 September 2026}
\begin{document}\maketitle

\section{Scope, source, and original objects}
This proof derives the gluing and collision assertions (32)--(46) of the
received account \cite{Received}, retaining the original odd generator
$\sigma$. The coefficient change to $\zeta$ below is accompanied by its full
inverse, residue cofactor, connection residue, and both original physical
receivers. The original ES labels are not merged when their coordinates
coincide. The marked prime is denoted $P$; the polynomial variable local to
it is $x=T-P$; the collision parameter is $t$. None is the fixed conductor
amplification $a_{\rm amp}$. The inherited polynomial provenance is the
Alp\"oge--Fable example as recorded by Tao \cite{Tao}; the four formal
$\mathrm M,\mathrm N,\mathrm T,\mathrm E$ labels and unchanged coordinate
marking are from the specified ES source \cite{ES}. The Gamma orthogonality
is the human source \cite{Gamma}; the original programme's FC20--24,
RS21--25 and GD14 specify its exact use \cite{Programme}. Classical
residue/dual-form provenance is \cite[Section 2]{JouveRV}. All finite
algebra and metric formulas required here are proved below.

Let the coefficient base be a characteristic-zero domain with $A,P$ units,
and let
\begin{equation}
 g(T)=T^3-s_1T^2+wT-s_3,\quad w=4s_3/P,\quad
 F(T)=A(T-P)g(T),\quad R=g(P).
 \tag{ECC1}
\end{equation}
Work where $R$ is not identically zero. The signed algebra is
$\mathscr E=B[T,\sigma]/((T-P)g,\sigma^2-F')$, in the unchanged even-then-odd
basis $(1,T,T^2,T^3,\sigma,\sigma T,\sigma T^2,\sigma T^3)$.
Its quotients are
\begin{equation}
 \mathscr E_P=B[\sigma]/(\sigma^2-AR),\qquad
 \mathscr E_g=B[T,\sigma]/(g,\sigma^2-A(T-P)g').
 \tag{ECC2}
\end{equation}
Division by the monic equations gives ranks $8,2,6$ respectively.

\section{Full gluing and the inverse CRT map}
The exact maps, including their target at the divisor, are
\begin{equation}
 0\longrightarrow\mathscr E\xrightarrow{j}\mathscr E_P\oplus\mathscr E_g
 \xrightarrow{d} B/(R)[\sigma]/(\sigma^2)\longrightarrow0.
 \tag{ECC3}
\end{equation}
Here $j$ restricts both coefficient polynomials and $d$ subtracts their
restrictions at $T=P$ modulo $R$. To prove injectivity, a polynomial
divisible by both $g$ and $T-P$ has form $gq$ with $R q(P)=0$.
The nonzerodivisor assumption on $R$ gives $q(P)=0$, hence divisibility by
$(T-P)g$. This applies to each even and odd coefficient. Their common
quotient has $(T-P,g)=(T-P,R)$ and $F'=AR=0$ there, giving the displayed
overlap including its odd nilpotent. For a compatible pair $(u,b(T))$,
divide $u-b(P)$ by $R$ and add that multiple of $g$ to $b$; this gives its
preimage. Every overlap element lifts from the prime factor, proving
surjectivity. This proves exactness and the fibre-product multiplication.

In the coefficient order $(u,b_0,b_1,b_2)$ the even comparison is
\begin{equation}
 \begin{gathered}
 C_R=\begin{pmatrix}1&P&P^2&P^3\\1&0&0&s_3\\
 0&1&0&-w\\0&0&1&s_1\end{pmatrix},\quad
 \det C_R=-R,\\
 C_R^{-1}=\begin{pmatrix}0&1&0&0\\0&0&1&0\\0&0&0&1\\0&0&0&0\end{pmatrix}
 +\frac1R\begin{pmatrix}-s_3\\w\\-s_1\\1\end{pmatrix}
             \begin{pmatrix}1&-1&-P&-P^2\end{pmatrix}.
 \end{gathered}
 \tag{ECC4}
\end{equation}
Indeed the inverse polynomial is
$b(T)+[u-b(P)]g(T)/R$. Evaluation and reduction prove both matrix products
equal $I_4$. The full signed matrix is $\diag(C_R,C_R)$, with determinant
$R^2$, full inverse $\diag(C_R^{-1},C_R^{-1})$, and two copies of the
displayed overlap. No sign state was deleted.

Put $\Delta_g=\operatorname{disc}(g)$. Then
\begin{equation}
 \det\Tr_{\mathscr E}=256A^4R^6\Delta_g^3,\qquad
 \det\Tr_{\mathscr E_P\oplus\mathscr E_g}=256A^4R^2\Delta_g^3.
 \tag{ECC5}
\end{equation}
For a monic quartic $f$, the ordinary even trace matrix has determinant
$\operatorname{disc}(f)$. In the signed algebra the trace matrix is block
diagonal with blocks $2T_f$ and $2T_fM_{Af'}$. Its determinant is
$2^8\operatorname{disc}(f)^2\prod_i Af'(r_i)
=256 A^4\operatorname{disc}(f)^3$, first on distinct roots and then
identically by polynomial continuation. Since
$\operatorname{disc}((T-P)g)=R^2\Delta_g$, the first formula follows.
The trace matrices under the inclusion differ by $C_R^{\oplus2}$ on
both sides, proving the second. Here ``trace'' is multiplication trace,
including on nonreduced fibres. These determinant identities by themselves
do not prove normality for an arbitrary base; the local normality statement
below includes a regular one-parameter base and a transverse collision.

\section{Original signed collision, exact forward and inverse maps}
Take a sufficiently small complex disc with coordinate $t$, and retain
analytic $P,A$ and the two remaining roots, such that $A$ never vanishes
and the latter roots remain distinct from one another, $P$, and $P+t$.
All formulas allow these units to depend analytically on $t$.
Write $h(P+x)=h_0+h_1x+x^2$, $h_t=h(P+t)$ and
\begin{equation}
 F=A x(x-t)h(P+x),\quad x^2=tx,\quad
 \sigma^2=A\{x(x-t)h'(P+x)+(2x-t)h(P+x)\}
          =A(2x-t)h(P+x).
 \tag{ECC6}
\end{equation}
The equality on the right is only in the displayed root quotient;
the full derivative on its left remains the original derivative.
Choose a single analytic square root of the nonzero function $Ah$ near
this disc and collision. Denote its values by
$q_0^2=Ah_0$, $q_t^2=Ah_t$, with $q_t=q_0$ at $t=0$.
The two values cannot be assigned unrelated signs if an analytic return
through $t=0$ is intended. The common sign remains a retained choice.
The unit in the root quotient and its inverse are exactly
\begin{equation}
 q(x)=q_0+\frac{q_t-q_0}{t}x,\qquad
 q(x)^{-1}=q_0^{-1}+\frac{q_t^{-1}-q_0^{-1}}t x.
 \tag{ECC7}
\end{equation}
The divided differences extend analytically. Multiplication followed by
$x^2=tx$ proves the inverse. Squaring and evaluating at $0,t$ proves
$q(x)^2=Ah(P+x)$; the same identity at zero follows by analytic continuity.

Define $\zeta=q(x)^{-1}\sigma$. Both directions of the isomorphism are
\begin{equation}
 \zeta^4=t^2,\quad x=(\zeta^2+t)/2,\quad
 \sigma=a\zeta+b\zeta^3,\quad
 a=(q_0+q_t)/2,\quad b=(q_t-q_0)/(2t),\quad
 \zeta=q(x)^{-1}\sigma.
 \tag{ECC8}
\end{equation}
Indeed $\zeta^2=2x-t$ and $x(x-t)=0$ are equivalent to the first two
relations. Substituting $q(x)$ proves the formula for $\sigma$, and
the inverse returns $x$ and $\sigma$ exactly. In the original basis
$e=(1,x,\sigma,x\sigma)$ and $z=(1,\zeta,\zeta^2,\zeta^3)$, let columns
transform by $[f]_z=S[f]_e$. Then
\begin{equation}
 S=\begin{pmatrix}1&t/2&0&0\\0&0&a&tq_t/2\\
 0&1/2&0&0\\0&0&b&q_t/2\end{pmatrix},\quad
 S^{-1}=\begin{pmatrix}
 1&0&-t&0\\0&0&2&0\\
 0&q_0^{-1}&0&-tq_0^{-1}\\
 0&(q_t^{-1}-q_0^{-1})/t&0&q_t^{-1}+q_0^{-1}
 \end{pmatrix},\quad \det S=-q_0q_t/4.
 \tag{ECC9}
\end{equation}
The last column of $S$ uses $x\sigma=(tq_t\zeta+q_t\zeta^3)/2$.
Each column of the inverse is ECC7 and
$\zeta^3=(2x-t)q(x)^{-1}\sigma$; this proves both inverse products and
invertibility at $t=0$. In particular, with $q=q_0|_{t=0}$ and
$\eta=h_1/h_0|_{t=0}$,
\begin{equation}
 \sigma=q\zeta+(q\eta/4)\zeta^3\quad(t=0).
 \tag{ECC10}
\end{equation}
Deleting the cubic term does not preserve the original odd coordinate.

The exact embedding of this local block in the full quartic algebra is
also retained. Put
$k_1=(h_t^{-1}-h_0^{-1})/t=-(h_1+t)/(h_0h_t)$ and
\begin{equation}
 e_c=h(P+x)(h_0^{-1}+k_1x),\quad e_h=1-e_c,\quad
 I_t=\begin{pmatrix}
 1&0\\h_1/h_0+k_1h_0&h_0/h_t\\
 1/h_0+k_1h_1&h_1/h_t\\k_1&1/h_t
 \end{pmatrix}.
 \tag{ECC11}
\end{equation}
Modulo $x(x-t)$, $e_c=1$; modulo $h$, it is zero. Hence $e_c^2=e_c$,
$e_h^2=e_h$, and $e_ce_h=0$ modulo the full quartic. The first column of
$I_t$ is $e_c$ and the second is $xe_c=xh/h_t$ modulo that quartic.
These claims follow either by division or by multiplication with
$x^4+(h_1-t)x^3+(h_0-th_1)x^2-th_0x=0$.
Consequently $\diag(I_t,I_t)S^{-1}$ is the entire four-to-eight original
coefficient embedding in $x$ coordinates, and translation $x=T-P$ uses
the full cubic binomial matrix, not a change of the original centres.
The complementary block $e_h\mathscr E$ has rank four and is retained.

\section{Normalization, specialization, and original connection residue}
Over the analytic one-dimensional base $B=\C\{t\}$, the normalization of
$B[\zeta]/(\zeta^4-t^2)$ is the sum of
$B[\zeta]/(\zeta^2+t)$ and $B[\zeta]/(\zeta^2-t)$.
Each summand is a regular local ring, isomorphic to $\C\{\zeta\}$ with
$t=\mp\zeta^2$; it is therefore integrally closed. The product is finite,
contains the original algebra and equals its total quotient algebra
after $t$ is inverted. Every integral element restricts to an integral
element of each normal factor, which proves the assertion.
In the order $(1,\zeta)$ at $x=0$ followed by $(1,\zeta)$ at $x=t$,
\begin{equation}
 C_\zeta=\begin{pmatrix}1&0&-t&0\\0&1&0&-t\\
 1&0&t&0\\0&1&0&t\end{pmatrix},\quad\det C_\zeta=4t^2,
 \qquad
 C_\sigma=\begin{pmatrix}1&0&0&0\\0&0&1&0\\
 1&t&0&0\\0&0&1&t\end{pmatrix},\quad\det C_\sigma=-t^2,
 \tag{ECC12}
\end{equation}
where $C_\sigma$ uses the original $(1,\sigma)$ in each factor.
They obey $C_\zeta S=\diag(1,q_0,1,q_t)C_\sigma$.
Their inverses send $(u_0,v_0,u_t,v_t)$ respectively to
\begin{equation}
 \left((u_0+u_t)/2,(v_0+v_t)/2,(u_t-u_0)/(2t),(v_t-v_0)/(2t)\right),
 \quad
 \left(u_0,(u_t-u_0)/t,v_0,(v_t-v_0)/t\right).
 \tag{ECC13}
\end{equation}
Elementary row subtraction and division by the units $2,q_0,q_t$
give Smith factors $(1,1,t,t)$, with cokernel $(B/(t))^2$.
Tensoring the short exact sequence with $B/(t)$ gives
\begin{equation}
 0\to\C^2\to\C[\zeta]/(\zeta^4)
 \to\C[\zeta]/(\zeta^2)\oplus\C[\zeta]/(\zeta^2)
 \to\C^2\to0.
 \tag{ECC14}
\end{equation}
The middle arrow sends a polynomial to its two remainders modulo
$\zeta^2$. Its kernel is precisely $(\zeta^2,\zeta^3)$; hence the initial
$\operatorname{Tor}_1$ injection has exactly this image. Its cokernel is
the difference of the two constant-linear coefficient pairs. Thus all
four original directions and the two lost specialization directions
are explicitly accounted for.

On the punctured disc differentiate the four labelled evaluation
branches. Since $t\partial_t\zeta^j=(j/2)\zeta^j$, the connection matrix
in $z$ is $D/t$, $D=\diag(0,1/2,1,3/2)$. In original coordinates it is
\begin{equation}
 \Omega_e=S^{-1}(DS/t+S'),\qquad
 \Res_{t=0}\Omega_e=
 \begin{pmatrix}0&0&0&0\\0&1&0&0\\
 0&0&1/2&0\\0&0&\eta/2&3/2\end{pmatrix}.
 \tag{ECC15}
\end{equation}
This is the derivative of $zS$ column by column. The residue is
$S(0)^{-1}DS(0)$ because $S'$ is analytic. Substitution of ECC9 proves
every entry, including the off-diagonal $\eta/2$. The normalized extension
has residue $\diag(0,1/2,0,1/2)$, so the traces differ by $2$, equal to
$t\partial_t\log\det C_\zeta$. Original $\sigma$ in its two factors
only adds analytic logarithmic derivatives of $q_0,q_t$, leaving these
residue eigenvalues unchanged. A loop in $t$ changes the signs over both
$P$ and $P+t$ without exchanging those labelled roots. This proves the
local monodromy assertion. ECC46--48 below supplies the further global
argument for the full normalized ES parameter space.

\section{Full original residue and trace inverse}
Use the original programme residue
$\Lambda(R+\sigma B)=\sum\Res B(T)dT/F(T)$, restricted to the two-root
block. For $B=b_0+b_1x$, at $t\ne0$ its exact values and continuation are
\begin{equation}
 \Lambda(\sigma B)=\frac{B(t)}{tq_t^2}-\frac{B(0)}{tq_0^2},\qquad
 \Lambda(\sigma)=\frac{q_t^{-2}-q_0^{-2}}t,\quad
 \Lambda(x\sigma)=q_t^{-2},\quad
 \Lambda(x^2\sigma)=tq_t^{-2}.
 \tag{ECC16}
\end{equation}
The formula follows from the two simple-root residues and extends by
the analytic divided difference; at zero it equals the double-root
residue $[x]B/(Ah(P+x))$. The full pairing in $e$ is
\begin{equation}
 \begin{gathered}
 B_e=\begin{pmatrix}0&J_e\\J_e&0\end{pmatrix},\quad
 J_e=\begin{pmatrix}(q_t^{-2}-q_0^{-2})/t&q_t^{-2}\\
 q_t^{-2}&tq_t^{-2}\end{pmatrix},\\
 \det J_e=-1/(q_0^2q_t^2),\quad
 J_e^{-1}=\begin{pmatrix}-tq_0^2&q_0^2\\
 q_0^2&(q_t^2-q_0^2)/t\end{pmatrix}.
 \end{gathered}
 \tag{ECC17}
\end{equation}
Odd-odd products are even and have zero residue; even-even products
also have zero residue. ECC16 proves the cross block, and direct
multiplication proves the displayed inverse. The residue pairing remains
perfect at the collision.

In $z$ the only nonzero values on its basis are
\begin{equation}
 \ell_1=\Lambda(\zeta)=\frac{q_t^{-3}-q_0^{-3}}t,\qquad
 \ell_3=\Lambda(\zeta^3)=q_t^{-3}+q_0^{-3}.
 \tag{ECC18}
\end{equation}
Indeed $\zeta=q^{-1}\sigma$, while $\zeta^3=(2x-t)q^{-1}\sigma$;
insert these two expressions into ECC16. With the monic transfer
functional $\lambda_z(f)=[\zeta^3](f\bmod(\zeta^4-t^2))$, the precise
residue cofactor and its inverse are
\begin{equation}
 \Lambda(f)=\lambda_z(cf),\quad
 c=\ell_3+\ell_1\zeta^2,\quad
 c^{-1}=\frac{q_0^3q_t^3}{4}(\ell_3-\ell_1\zeta^2),\quad
 c|_{t=0}=\frac2{q^3}\left(1-\frac{3\eta}{4}\zeta^2\right).
 \tag{ECC19}
\end{equation}
Here $\ell_3^2-t^2\ell_1^2=4/(q_0^3q_t^3)$ proves the inverse. Both
functionals vanish on even powers, so testing the two odd basis powers
proves the functional identity completely. The limit uses
$\partial_x q/q=h_1/(2h_0)$. The full Gram in $z$ is
\begin{equation}
 B_z=\begin{pmatrix}0&\ell_1&0&\ell_3\\
 \ell_1&0&\ell_3&0\\0&\ell_3&0&t^2\ell_1\\
 \ell_3&0&t^2\ell_1&0\end{pmatrix},\quad
 \det B_z=16/(q_0^6q_t^6),\quad B_e=S^{\mathsf T}B_zS.
 \tag{ECC20}
\end{equation}
This also checks the constant and sign in $\det S$ against ECC17.

Put $a_0=q_0^2$, $a_t=q_t^2$, and use the original even order $(1,x)$.
The full odd multiplication, the trace multiplier, and their inverses are
\begin{equation}
 \begin{gathered}
 H=\begin{pmatrix}-ta_0&0\\a_0+a_t&ta_t\end{pmatrix},\qquad
 M_\sigma=\begin{pmatrix}0&H\\I_2&0\end{pmatrix},\quad
 M_\sigma^{-1}=\begin{pmatrix}0&I_2\\H^{-1}&0\end{pmatrix},\\
 H^{-1}=\begin{pmatrix}-1/(ta_0)&0\\
 (a_0+a_t)/(t^2a_0a_t)&1/(ta_t)\end{pmatrix},\\
 H^{-2}=\begin{pmatrix}1/(t^2a_0^2)&0\\
 (a_0+a_t)(a_0-a_t)/(t^3a_0^2a_t^2)&1/(t^2a_t^2)\end{pmatrix},\\
 \Theta=M_{2\sigma^3}=2\begin{pmatrix}0&H^2\\H&0\end{pmatrix},\qquad
 \Theta^{-1}=\frac12\begin{pmatrix}0&H^{-1}\\H^{-2}&0\end{pmatrix},\quad
 \det\Theta=-16t^6a_0^3a_t^3.
 \end{gathered}\tag{ECC21}
\end{equation}
The original relation ECC6 gives $\sigma^2=-ta_0+(a_0+a_t)x$,
which proves $H$. The inverses follow by triangular multiplication;
the determinant follows from $\det H=-t^2a_0a_t$ and the even block swap.
Multiplication trace obeys $\Tr M_f=\Lambda(2\sigma^3f)$: it is the sum
of all four values when $t\ne0$, and both sides extend analytically.
Thus $B_e\Theta$ is the full trace pairing, not the perfect residue
pairing. At $t=0$, with $a=q^2$, it has rank one, while $B_e$ has rank four.

\section{Both original Gamma receivers, with their full finite matrices}
For any scalar $c$ define
\begin{equation}
 C_c=\begin{pmatrix}1&c&c^2&c^3\\0&1&2c&3c^2\\
 0&0&1&3c\\0&0&0&1\end{pmatrix},\quad
 [f(r+c)]=C_c[f(r)],\quad C_c^{-1}=C_{-c}.
 \tag{ECC22}
\end{equation}
Binomial expansion proves all entries and both inverse products.
Hold the original conductor, its period, actual order $v$, and moments
$\mu_v,\ldots,\mu_{v+3}$ fixed, with $\mu_v\ne0$. Put
$c_*=a_{\rm amp}/2-4$, $\kappa=c_*-1/2$, $b=s/2>0$ and
$M=(2\pi)^{s/2}$ as in the original Gamma measure. Its actual source and
target are $U=\mathcal P_{v+3}/\mathcal P_{v-1}$ and $W=\mathcal P_3$.
Their conductor is the given $T_A$ (the subscript here is the original
conductor notation, not multiplication by the quartic coefficient).
The scalar halves of GD14, on coefficient columns in $r$, are
\begin{equation}
 F_W f=f(S'-\kappa),\quad F_U f=T_A^{-1}f(S'-\kappa),\quad
 F_W=T_AF_U,\quad F_W^{-1}w=w(r+\kappa),\quad
 F_U^{-1}u=(T_Au)(r+\kappa).
 \tag{ECC23}
\end{equation}
Their signed receivers are $F_W^{\oplus2},F_U^{\oplus2}$, not an
identification of multiplication by $r$ with the original conductor.

Here are finite explicit original Gram matrices and their inverses.
Use the unchanged shifted moments and phases
\begin{equation}
 \begin{gathered}
 \nu_{v+j}=\sum_{k=0}^{j}\binom{v+j}{v+k}(-4)^{j-k}\mu_{v+k},\\
 g_0=(-i)^v\mu_v/v!,\quad g_1=(-i)^{v+1}\nu_{v+1}/(v+1)!,\\
 g_2=(-i)^{v+2}\nu_{v+2}/(v+2)!-(v/3)g_0,\quad
 g_3=(-i)^{v+3}\nu_{v+3}/(v+3)!-((v+1)/3)g_1,\\
 \xi_0=g_0^{-1},\quad\xi_1=-g_1/g_0^2,\quad
 \xi_2=g_1^2/g_0^3-g_2/g_0^2,\\
 \xi_3=-g_1^3/g_0^4+2g_1g_2/g_0^3-g_3/g_0^2,\\
 k_j=(b)_{v+j}/(v+j)!,\quad D_i=\diag(1,i,-1,-i).
 \end{gathered}\tag{ECC24}
\end{equation}
The cubic products $(\sum g_jz^j)(\sum\xi_jz^j)=1\bmod z^4$
prove every inverse coefficient; the shifted moment formula is the
coefficient expansion of $e^{-4z}E_A(z)$. Define
\begin{equation}
 \begin{gathered}
 V=\begin{pmatrix}\sqrt M&0&b\sqrt M&0\\
 0&\sqrt{Mb}&0&(3b+2)\sqrt{Mb}\\
 0&0&\sqrt{2Mb(b+1)}&0\\0&0&0&\sqrt{6Mb(b+1)(b+2)}\end{pmatrix},\\
 F=\begin{pmatrix}\xi_0&\xi_1&b\xi_0+2\xi_2&(3b+2)\xi_1+6\xi_3\\
 0&\xi_0&2\xi_1&(3b+2)\xi_0+6\xi_2\\
 0&0&2\xi_0&6\xi_1\\0&0&0&6\xi_0\end{pmatrix},\\
 B_F=\begin{pmatrix}g_0&g_1&g_2-bg_0/2&g_3-bg_1/2\\
 0&g_0&g_1&g_2-(3b+2)g_0/6\\0&0&g_0/2&g_1/2\\0&0&0&g_0/6\end{pmatrix},\\
 W_x=VD_i,\qquad U_x=\sqrt M\diag(\sqrt{k_j})FD_i,\qquad
 U_x^{-1}=M^{-1/2}D_i^*B_F\diag(k_j^{-1/2}).
 \end{gathered}\tag{ECC25}
\end{equation}
The target orthonormal polynomials are $p_j((S'-c_*)/i)/\sqrt{Mj!(b)_j}$,
where $p_0=1,p_1=y,p_2=y^2-b,p_3=y^3-(3b+2)y$.
Expanding their monomials gives $V$, and their source inverse images
under FC24 give $U_x$. More explicitly that inverse in orthonormal
coordinates has entries $\xi_{j-r}\rho_j/\rho_{v+r}$ for $r\le j$, where
$\rho_j=\sqrt{j!/(b)_j}$; multiplying $VD_i$ gives the displayed $U_x$.
The convolution identity in ECC24 proves $FB_F=B_FF=I_4$.
This supplies a direct finite check, with all original masses present.

In $r$ coefficients the actual two metrics and complete covariance
factors are
\begin{equation}
 \begin{gathered}
 H_U=C_{1/2}^*U_x^*U_xC_{1/2},\quad
 H_W=C_{1/2}^*W_x^*W_xC_{1/2},\\
 Q_U=C_{-1/2}U_x^{-1},\quad Q_W=C_{-1/2}D_i^*V^{-1},\quad
 H_X^{-1}=Q_XQ_X^*\quad(X=U,W).
 \end{gathered}\tag{ECC26}
\end{equation}
For specificity $V^{-1}$ has diagonal
$M^{-1/2},(Mb)^{-1/2},[2Mb(b+1)]^{-1/2},[6Mb(b+1)(b+2)]^{-1/2}$,
entry $(0,2)$ equal $-b[2Mb(b+1)]^{-1/2}$, entry $(1,3)$ equal
$-(3b+2)[6Mb(b+1)(b+2)]^{-1/2}$, and all other off-diagonal entries zero.
Direct multiplication proves this inverse. ECC26 follows because
$r=S'-\kappa=(S'-c_*)+1/2$. Every displayed matrix is invertible, so
both metrics are strictly positive. No original moment has been assigned
a substitute value.

The received coefficient receiver $\Psi_*$ is different from literal
substitution. Retain its exact definition
$\Phi_*=(V_*^{\mathsf T})^{-1}K^{-1}$ and
$\Psi_*=B_{v,*}(V_*^{\mathsf T})^{-1}K^{-1}=T_A\Phi_*$ in the inherited
coordinate bases. Its exact correction on $r$ columns is
\begin{equation}
 M_* = C_\kappa\Psi_*,\quad
 M_*^{-1}=K V_*^{\mathsf T} B_{v,*}^{-1}C_{-\kappa},\quad
 F_WM_* = \Psi_*,\quad F_UM_* = \Phi_*,\quad
 H_X^{\rm rec}=M_*^*H_XM_*.
 \tag{ECC27}
\end{equation}
The fixed matrices $K,V_*,B_{v,*}$ here are precisely the original
ES marking, reference evaluation and conductor matrices of FC26--32,
not matrices recomputed from the colliding family. Multiplying ECC23
by $M_*$ proves each identity; their inverses prove equivalence of the
two receivers. Consequently the received covariance is obtained by
replacing $Q_X$ with $M_*^{-1}Q_X$. On physical spaces the exact maps
$L_W=\Psi_*F_W^{-1}$ and $L_U=\Phi_*F_U^{-1}$ obey
$T_AL_U=L_WT_A$ by multiplication. All results below apply separately
to $H_U,H_W$ and, with this correction, their received versions.

\section{The complete minimum and all raw signed eigenvalues}
Fix one of the positive matrices $H=H_X$ from ECC26, or the explicitly
corrected received matrix. Let $v(r)=(1,r,r^2,r^3)^{\mathsf T}$ and
\begin{equation}
 q_j(r)=v(r)^{\mathsf T}Qe_j,\qquad
 K(r,s)=\sum_{j=0}^3q_j(r)\overline{q_j(s)}
       =v(r)^{\mathsf T}H^{-1}\overline{v(s)}.
 \tag{ECC28}
\end{equation}
Here $Q$ is explicitly ECC26, or $M_*^{-1}Q_X$. Thus this is a finite
formula in the original moments and masses, not an unspecified inverse.
The signed norm of $(f_0,f_1)$ is $[f_0]^*H[f_0]+[f_1]^*H[f_1]$.
For the observation $(f_0(P),f_1(P))=(u,v)$ the complete minimum is
\begin{equation}
 \frac{|u|^2+|v|^2}{K(P,P)},\quad
 f_0(r)=uK(r,P)/K(P,P),\quad f_1(r)=vK(r,P)/K(P,P).
 \tag{ECC29}
\end{equation}
Indeed $H^{-1}\overline{v(P)}$ is the representing vector of evaluation;
its squared norm is $K(P,P)>0$. Every other representative is this
one plus a vector in the entire evaluation kernel, orthogonal to it.
Pythagoras proves attainment, uniqueness, and the formula separately
for each coefficient. It includes the complementary quartic directions.
In particular the CRT section $g/R$ has the exact excess
\begin{equation}
 \|g/R\|_H^2=\frac1{K(P,P)}+
 \left\|g/R-K(\,\cdot,P)/K(P,P)\right\|_H^2,\qquad
 \frac{\det G_{\rm CRT}}{\det G_{\min}}=
 \left(\frac{K(P,P)\|g\|_H^2}{|R|^2}\right)^2.
 \tag{ECC30}
\end{equation}
The two signed section columns occupy the separate coefficient copies,
so their Gram is $\|g/R\|_H^2I_2$ and the minimum Gram is
$K(P,P)^{-1}I_2$. This proves the determinant ratio with no removal of
kernel directions. The minimum stays positive at the collision although
the CRT section has the displayed pole.

For $t\ne0$ let $s_0^2=-tq_0^2$, $s_t^2=tq_t^2$, retaining both choices
of $s_0,s_t$. Order the four raw values as $(P,+),(P,-),(P+t,+),(P+t,-)$.
Their observation matrix on the eight original coefficients is
\begin{equation}
 O_t=\begin{pmatrix}
 v(P)^{\mathsf T}&s_0v(P)^{\mathsf T}\\
 v(P)^{\mathsf T}&-s_0v(P)^{\mathsf T}\\
 v(P+t)^{\mathsf T}&s_tv(P+t)^{\mathsf T}\\
 v(P+t)^{\mathsf T}&-s_tv(P+t)^{\mathsf T}
 \end{pmatrix},\quad
 C_t=O_t\diag(H^{-1},H^{-1})O_t^*.
 \tag{ECC31}
\end{equation}
It has rank four: invert each signed sum/difference and then interpolate
the two distinct root values separately. A unitary signed sum/difference,
followed by ordering the two even observations first, gives the full
covariance blocks $2K_2$ and $2D_sK_2D_s^*$, where
\begin{equation}
 K_2=\begin{pmatrix}K_0&C\\\overline C&K_t\end{pmatrix},\quad
 K_0=K(P,P),\ K_t=K(P+t,P+t),\ C=K(P,P+t),\quad
 D_s=\diag(s_0,s_t).
 \tag{ECC32}
\end{equation}
Consequently all four finite covariance eigenvalues, with no asymptotic
replacement, are the sorted union
\begin{equation}
 \begin{aligned}
 \lambda_{E,\pm}&=K_0+K_t\ \pm\sqrt{(K_0-K_t)^2+4|C|^2},\\
 \lambda_{O,\pm}&=a+b\ \pm\sqrt{(a-b)^2+4|s_0s_t|^2|C|^2},\qquad
 a=|s_0|^2K_0,\quad b=|s_t|^2K_t.
 \end{aligned}\tag{ECC33}
\end{equation}
These follow from the quadratic characteristic polynomials of the two
Hermitian blocks. Positivity follows from rank two of the two evaluation
rows and positivity of $H$. In particular
$\lambda_{E,+}\lambda_{E,-}=4\det K_2$ and
$\lambda_{O,+}\lambda_{O,-}=4|s_0s_t|^2\det K_2$.

Take $P$ fixed in the limit, or use its analytic limiting value. Write
\begin{equation}
 K=\sum_j|q_j(P)|^2,\quad
 J=\begin{pmatrix}
 K&\sum_jq_j(P)\overline{q_j'(P)}\\
 \sum_jq_j'(P)\overline{q_j(P)}&\sum_j|q_j'(P)|^2
 \end{pmatrix},\quad L=\det J/K>0,\quad H_0=|Ah_0|_{t=0}.
 \tag{ECC34}
\end{equation}
The two rows $v(P)^{\mathsf T}Q$ and $v'(P)^{\mathsf T}Q$ are independent,
because $Q$ is invertible and evaluation/derivative are independent on
cubics. Thus $J$ is positive definite. Expanding its determinant gives
$KL=\sum_{i<j}|q_i(P)q_j'(P)-q_j(P)q_i'(P)|^2$, proving the exact
positive constant. Divided differences give
$\det K_2=|t|^2(KL+o(1))$, while $K_t,C\to K$ and
$|s_0|^2/|t|,|s_t|^2/|t|\to H_0$. Therefore, in decreasing order,
\begin{equation}
 \lambda_1\to4K,\quad
 \lambda_2/|t|\to4H_0K,\quad
 \lambda_3/|t|^2\to L,\quad
 \lambda_4/|t|^3\to H_0L.
 \tag{ECC35}
\end{equation}
For the smaller eigenvalue of each block use its exact determinant
divided by its larger eigenvalue in ECC33. This proves all four limits
also for arbitrary complex approach to zero; no phase is suppressed in
the exact matrices. The four different powers force the stated order
for all sufficiently small nonzero $t$.

The attained right inverse, its metric, and every exterior cost are
\begin{equation}
 \begin{gathered}
 R_t=\diag(H^{-1},H^{-1})O_t^*C_t^{-1},\quad
 O_tR_t=I_4,\\ R_t^*\diag(H,H)R_t=C_t^{-1},\quad
 \|\wedge^jR_t\|=\prod_{k=5-j}^4\lambda_k^{-1/2}\quad(1\le j\le4).
 \end{gathered}
 \tag{ECC36}
\end{equation}
Every other solution of the raw values is $R_tu+\ker O_t$, orthogonal
in the original metric, as multiplication of the displayed matrices
shows. Thus this is the full original quotient inverse, not the inverse
of a selected algebraic section. Its singular values are the square
roots of the eigenvalues of its Gram $C_t^{-1}$; singular-value
decomposition proves the exterior product formula. In particular
\begin{equation}
 \begin{aligned}
 |t|^{3/2}\|R_t\|&\longrightarrow(H_0L)^{-1/2},\\
 |t|^{5/2}\|\wedge^2R_t\|&\longrightarrow(H_0^{1/2}L)^{-1},\\
 |t|^3\|\wedge^3R_t\|&\longrightarrow(2H_0L\sqrt K)^{-1},\\
 |t|^3\|\wedge^4R_t\|&\longrightarrow(4H_0KL)^{-1},\qquad
 \|\wedge^4R_t\|=\frac1{4|s_0s_t|\det K_2}.
 \end{aligned}\tag{ECC37}
\end{equation}
The determinant expression follows directly from the two covariance
blocks. ECC33 and ECC36 are exact at every admissible nonzero parameter,
so the bounds are not limited to these leading terms.

\section{Retained jet map and original quotient trace cost}
The exact raw-to-jet matrix, preserving original $\sigma$ values, is
\begin{equation}
 Z_t=\begin{pmatrix}
 1/2&1/2&0&0\\-1/(2t)&-1/(2t)&1/(2t)&1/(2t)\\
 1/(2s_0)&-1/(2s_0)&0&0\\
 -1/(2ts_0)&1/(2ts_0)&1/(2ts_t)&-1/(2ts_t)
 \end{pmatrix},\quad
 |\det Z_t|=\frac1{4|s_0s_tt^2|}.
 \tag{ECC38}
\end{equation}
It sends the four values to
$(f_0(P),(f_0(P+t)-f_0(P))/t,f_1(P),(f_1(P+t)-f_1(P))/t)$.
Its inverse sends $(a,b,c,d)$ to
$(a+s_0c,a-s_0c,a+tb+s_t(c+td),a+tb-s_t(c+td))$.
These four identities prove both inverse products and the determinant.
Let
\begin{equation}
 E_t=\begin{pmatrix}v(P)^{\mathsf T}\\
 [v(P+t)^{\mathsf T}-v(P)^{\mathsf T}]/t\end{pmatrix},\quad
 J_t=E_tH^{-1}E_t^*,\quad G_t=J_t^{-1}.
 \tag{ECC39}
\end{equation}
Then $Z_tO_t=\diag(E_t,E_t)$ and
$Z_tC_tZ_t^*=\diag(J_t,J_t)$. At zero, $E_0$ has rows
$v(P)^{\mathsf T},v'(P)^{\mathsf T}$, so $J_t\to J>0$ and
$G_t\to G=J^{-1}$. The full minimum metric on the original local
coefficient algebra $(1,x,\sigma,x\sigma)$ is
$\diag(G_t,G_t)$. This follows from the same representing-vector
calculation as ECC36 with $E_t$ in place of $O_t$ and includes the entire
kernel of the quartic-to-local quotient. It is different from assigning
a Euclidean metric to local coordinates or restricting to the idempotent
section ECC11. Explicitly $G_t$ is the two-by-two adjugate of $J_t$
divided by $\det J_t>0$.

ECC21 now has a complete original quotient-metric inverse cost. For
$F=H^{-1},H^{-2}$ put
\begin{equation}
 a_F=\Tr(G_t^{-1}F^*G_tF),\quad b_F=|\det F|^2,\quad
 \lambda_{F,\pm}=(a_F\pm\sqrt{a_F^2-4b_F})/2.
 \tag{ECC40}
\end{equation}
The four singular values of $\Theta^{-1}$ in
$\diag(G_t,G_t)$ are the sorted union
$\tfrac12\sqrt{\lambda_{H^{-1},\pm}}$ and
$\tfrac12\sqrt{\lambda_{H^{-2},\pm}}$. Indeed the two off-diagonal
blocks act between orthogonal identical metric copies, and the
characteristic polynomial of each positive squared singular matrix is
$z^2-a_Fz+b_F$. This proves exact finite values and all exterior norms
by products, with no discarded direction.

Here are also all leading constants. Put $a=(Ah_0)|_{t=0}$,
$\eta=(h_1/h_0)|_{t=0}$, $N=\left(\begin{smallmatrix}0&0\\1&0\end{smallmatrix}\right)$,
and $n^2=G_{11}(G^{-1})_{00}=K/L$ (indices are zero-based). ECC21 gives
\begin{equation}
 \lim_{t\to0}t^2\Theta^{-1}=
 \frac12\begin{pmatrix}0&2N/a\\a^{-2}(I_2-2\eta N)&0\end{pmatrix},\qquad
 \Theta(0)=\begin{pmatrix}0&0\\4aN&0\end{pmatrix}.
 \tag{ECC41}
\end{equation}
The difference $a_t-a_0=A t(h_1+t)$ proves the limit even when $A,h_0$
depend on $t$. The first limit has rank three and its positive singular
values are the sorted list $(c_1,c_2,c_3)$ of
\begin{equation}
 \left\{\frac n{|a|},\quad
 \frac{\sqrt{1+|\eta|^2n^2}+|\eta|n}{2|a|^2},\quad
 \frac{\sqrt{1+|\eta|^2n^2}-|\eta|n}{2|a|^2}\right\}.
 \tag{ECC42}
\end{equation}
To verify it, the rank-one $N$ has squared metric norm
$\Tr(G^{-1}N^*GN)=G_{11}(G^{-1})_{00}=n^2$ and trace zero.
Consequently $I-2\eta N$ has determinant one and squared Frobenius norm
$2+4|\eta|^2n^2$, yielding the two displayed singular values by the
quadratic formula. The limiting forward matrix in ECC41 has its only
nonzero singular value $4|a|n$. Inverting singular values and using
continuity gives three inverse singular values $c_j|t|^{-2}(1+o(1))$
and the fourth tending to $(4|a|n)^{-1}$. Thus, with
$d=(c_1,c_2,c_3,(4|a|n)^{-1})$,
\begin{equation}
 \lim_{t\to0}|t|^{2\min(j,3)}\|\wedge^j\Theta^{-1}\|
 =\prod_{k=1}^j d_k\quad(1\le j\le4),\qquad
 \|\wedge^4\Theta^{-1}\|=\frac1{16|t|^6|a_0a_t|^3}.
 \tag{ECC43}
\end{equation}
The determinant norm equals the absolute determinant because the same
full metric is used in domain and codomain. The product of the constants
is $1/(16|a|^6)$, agreeing with ECC21. This is the trace-to-residue
multiplier in the actual local quotient; the original fixed $T_A$
has not been replaced by $\Theta$.

\section{An exact ES path and its original constants}
The explicit labelled algebraic ES family in the received source is
\begin{equation}
 (P;P+t,2P,z_t),\quad z_t=\frac{2P(P+t)}{5P+7t},\quad
 A=-\frac{5P+7t}{22P^2+35Pt+7t^2},\quad
 h_0=-\frac{P^2(3P+5t)}{5P+7t},\quad h_1=-z_t.
 \tag{ECC44}
\end{equation}
Direct addition gives
$1/(P+t)+1/(2P)+1/z_t=4/P$; adding all four roots gives the displayed
unchanged leading coefficient $A$. Also
$h(P+x)=(x-P)(x+P-z_t)$, proving $h_0,h_1$. Near zero all required
units are nonzero when $P\ne0$. At $t=0$ the exact constants are
\begin{equation}
 a=3P/22,\quad H_0=|3P/22|,\quad\eta=2/(3P),\quad
 \sigma=q\zeta+\frac q{6P}\zeta^3,\quad
 c=\frac2{q^3}\left(1-\frac{\zeta^2}{2P}\right),\quad q^2=3P/22.
 \tag{ECC45}
\end{equation}
These specialize ECC10, ECC19 and every original-metric constant in
ECC35--43, retaining $K_U,L_U$ and $K_W,L_W$ separately. The local residue
shear in ECC15 is $1/(3P)$. The four signed branches have values
$\sigma=\pm q_0\sqrt{-t}$ and $\sigma=\pm q_t\sqrt t$, with their two
root labels unchanged; thus the minimum, covariance, jet and trace costs
all concern the same four states with the original odd coordinate.
This is a complex algebraic family, not an assertion of integer ES
witnesses at arbitrary $t$.

Every map into either original physical space is now explicit:
translate ECC11 from $x$ to $r$ with $C_{-P}$, apply the two ECC23
receivers, and use ECC9 for $\zeta$ columns when required. The source and
target maps intertwine by $T_A\oplus T_A$ identically, since ECC23 does.
For local quotients their respective minimum metrics and all costs are
ECC26--43. These observations prove the exact return and all its finite
costs. The nonproper signed chart, the degenerating trace, and the
raw-point observation have explicit losses; the fixed original conductor
itself remains invertible with its unchanged nonzero moment $\mu_v$.
No arithmetic order jump or Riemann-hypothesis conclusion is inferred
from a different map's pole.

\section{The complete ES signed monodromy}
Here use the actual normalized ES parameter space
$\operatorname{Spec}\C[A,A^{-1},P,P^{-1},W]$, with
$s_1=-A^{-1}-P$ and $s_3=PW/4$, removing the divisors where roots
coincide. The four roots keep their prime/denominator marking. Put
$c=A^{-1}+P$. Over the function field $\C(P,c,W)$ the denominator
polynomial is
\begin{equation}
 g(T)=T^2(T+c)+W(T-P/4).
 \tag{ECC46}
\end{equation}
It is irreducible. In fact the two coefficients of this polynomial as
a degree-one polynomial in $W$ are coprime in $\C(P,c)[T]$:
$T=P/4$ is not a root of $T^2(T+c)$. A factorization in the two-variable
polynomial ring would therefore have a factor of $W$-degree zero
dividing both coefficients, which is impossible unless it is a unit.
Gauss's lemma then proves irreducibility in $\C(P,c,W)[T]$. Thus the
permutation group on the three denominators is transitive.

There is an explicit denominator-transposition loop within this ES
base. Fix $P\ne0$ and $a=2P$ and use the parameter $u$ with roots
\begin{equation}
 P,\quad a+\sqrt u,\quad a-\sqrt u,\quad
 z(u)=\frac{P(a^2-u)}{4(a^2-u)-2Pa},\qquad
 A(u)=-\frac1{P+2a+z(u)}.
 \tag{ECC47}
\end{equation}
Their denominator reciprocal sum is $4/P$ by direct addition. All
quartic coefficients are analytic in $u$, because the two moving roots
enter symmetrically. At zero the latter denominator is $P/3$, distinct
from both $P$ and $a$. The denominator discriminant is
$4u[(a-z(u))^2-u]^2$, with nonzero coefficient of $u$ at zero. Therefore
a small loop in $u$ exchanges exactly the two coalescing denominators.
A transitive subgroup of $S_3$ containing a transposition is $S_3$:
its order is divisible by both $3$ and $2$, hence equals $6$.

At a base point label the four roots by $r_0=P,r_1,r_2,r_3$ and their
odd square roots by $\sigma_i^2=A\prod_{j\ne i}(r_i-r_j)$. The product
of these four equations is
$(\prod\sigma_i)^2=A^4\prod_{i<j}(r_j-r_i)^2$. The quotient
$\prod\sigma_i/[A^2\prod_{i<j}(r_j-r_i)]$ has values in $\{1,-1\}$
and is constant along each lifted path. Thus every monodromy element
$(\epsilon,\pi)$ satisfies
\begin{equation}
 \pi\in S_3\text{ fixes }0,\qquad
 \epsilon_i\in\{1,-1\},\qquad
 \prod_{i=0}^3\epsilon_i=\operatorname{sgn}\pi.
 \tag{ECC48}
\end{equation}
This is an upper bound of $8\cdot6=48$ elements. The prime--denominator
loop ECC44, with ECC8's analytic unit, changes precisely the two signs
at $P$ and $P+t$ and fixes the root permutation; the other two signed
roots stay in nonvanishing analytic discs and do not change sign.
All loops may be based in the same connected complement of the
discriminant by joining their initial points with paths; this complement
is connected because it is a nonempty Zariski-open subset of the
irreducible smooth complex parameter space. Conjugation by the already
proved denominator $S_3$ gives the three flips $(0,1),(0,2),(0,3)$.
They are independent vectors in the four-bit sign space and generate
its full eight-element even subgroup. Together with a surjection onto
$S_3$, this gives at least $48$ elements. The upper bound proves exact
equality with ECC48. Thus the global group assertion is proved on this
precisely specified ES base, in addition to the full local original
coordinate and metric maps above.

\begin{thebibliography}{9}
\bibitem{Received} Received AI-generated web calculation,
\emph{The ES bridge at the defining prime of an actual Erd\H{o}s--Straus witness},
20 September 2026, equations (32)--(46). No human author supplied.
Complete received SHA-256:
\texttt{f97ec18809279c10caf20a0ffa59abaa9cf54ec5c5c0ab21761c2840b426d886}.
The whole received file was read; the present proof independently derives
the displayed gluing, collision, monodromy and metric statements with
the stated corrections and the specified coefficient base.
\bibitem{Programme} Split-Zero programme, \emph{The ES--Fable inverse
correspondence in the original weighted conductor}, cumulative source,
20 September 2026, FC20--24, FC26--32, RD1--12, RS21--25 and GD14--17.
The complete original source and its cited human provenance are retained
with the receiving programme edition.
\bibitem{ES} The Clankers, \emph{Erd\H{o}s--Straus project reader},
canonical 488-page edition, source version 69,
\texttt{erdos\_straus\_project\_reader.tex}, theorem
\texttt{explicit-four-time-Jacobian-collision}, lines 19484--19588.
Source SHA-256:
\texttt{678b3f70ee7d2a6d00f6c588d7778e390118ab4743a5759612dbba6430984c78}.
\bibitem{Tao} Terence Tao, \emph{A digestion of the Jacobian conjecture
counterexample}, 21 July 2026, equation (1), recording the Alp\"oge--Fable
example. \url{https://terrytao.wordpress.com/2026/07/21/a-digestion-of-the-jacobian-conjecture-counterexample/}.
The formal ES marking is additional programme data, not a claim about
a fourth point in the three-dimensional human example.
\bibitem{Gamma} T. H. Koornwinder, R. Wong, R. Koekoek, R. F. Swarttouw
and W. P. Reinhardt, NIST Digital Library of Mathematical Functions,
Chapter 18, equations 18.19.8--9, 18.22.8 and 18.23.7.
\url{https://dlmf.nist.gov/18}.
Original orthogonality is inherited through WCF1--6 and RS21--25, whose
unchanged masses, phases and centres are included in ECC24--26.
\bibitem{JouveRV} Florent Jouve and Fernando Rodriguez-Villegas,
\emph{On the bilinear structure associated to Bezoutians}, Journal of
Algebra 400 (2014), 161--184, DOI 10.1016/j.jalgebra.2013.10.030,
original author TeX arXiv:1004.4208v3, Section 2, lemma \texttt{Q-lemma}
and theorem \texttt{thBsharp}. \url{https://arxiv.org/abs/1004.4208v3}.
These are the classical transfer-functional sources cited in RD;
ECC16--21 proves the present full signed calculation directly.
\end{thebibliography}
\end{document}
